\SourceLeaf{338}%
\par\indent Für die erste Näherung setzt man \(T=5^0\ 42'\ 22''\) und liest in der Tafel zum Argument \(51^0\ 48'\) den logarithmischen Wert \((2)=8{,}51004\) ab. So erhält man \(\beta=3412'',\ B=51^0 19'36''\). Setzt man damit \(\log(1)=8{,}50893\), so erhält man \(\tau=425\) und \(T=5^0\ 38'\ 49''\). Eine neue Rechnung mit diesem Wert, bei der man mit dem zuvor gefundenen Wert von \(B\) nun \(\log(2)=8{,}51007\) setzt, ergibt
\[
\beta=3413'',\quad B=51^0\ 19'\ 35'',\quad\tau=420''\ 5,\quad T=5^0\ 38'\ 51''\ 5
\]
Mit dem gefundenen Wert von \(B\) entnimmt man der Tafel
\[
\begin{aligned}
\log(1)&=8{,}5089337\\
\log(2)&=8{,}5100716\\
\log(3)&=1{,}94876\\
\log(4)&=3{,}32553\\
\log(5)&=4{,}92770\\
\log(6)&=4{,}61132
\end{aligned}
\]
Mit \(T=5^0\ 38'\ 51''\ 5\) erhält man zunächst \(\log\beta=3{,}5331341\) oder \(\beta=3412''\ 983\). Setzt man nochmals \(\beta\) statt \(b\) ein, erhält man \(B=51^0 19'35''\ 4379\). Damit ergibt sich ferner \(\log\tau=2{,}6238475\). Als Nächstes findet man in Einheiten der siebten Dezimalstelle
\[
\begin{aligned}
(3)rr&=99{,}80\\
(4)\beta\beta&=2{,}46\\
(5)\beta\beta&=98{,}62\\
(6)\beta\beta&=47{,}60\\
3(7)\tau\tau&=2{,}26\\
(7)\tau\tau&=0{,}75
\end{aligned}
\]
und damit die logarithmischen Korrekturen
\[
\begin{aligned}
\text{von}\quad\log\beta\quad .\ .\ .\ .\ .\quad&+\phantom{10}3\\
\log\tau\quad .\ .\ .\ .\ .\quad&+103\\
\log\lambda\quad .\ .\ .\ .\ .\quad&-\phantom{1}47
\end{aligned}
\]
Man hätte diese Rechnung auch schon mit den früheren Werten von \(\log\beta\) und \(\log\tau\) durchführen können und hätte dasselbe Ergebnis erhalten. Dann hätten sich mit \(\log b=3{,}5331344\) sofort die Werte \(b=3412''\ 985\) und \(B=51^0 19'35''4369\)
